\documentclass[11pt]{scrartcl}
\usepackage{ebb}



\title {Conjugados Isogonales}
\subtitle {Entrenamientos Nacionales Centros Pagmos}

\author{Emmanuel Buenrostro}

\begin{document}
\maketitle

\section{Teoria}

En construcción jeje.

\begin{itemize}
    \item Qué es isogonal 
    \item Simediana
    \item Existencia Conjugado isogonal (Ceva trigonometrico )
    \item Nueve puntos para isogonales
    \item 
    \item First isogonality lemma
\end{itemize}

\section{Problemas}
\epigraph{It is whispered that any child who mistakes DRIFLOON for a balloon and holds on to it could wind up missing.}{ }

\subsection{Nivel I}
\problemdiff{0}
\begin{problem}
    Demuestra que $H$ y $O$ son conjugados isogonales.
\end{problem}

\begin{problem}
    Demuestra que $I$ e $I$ son conjugados isogonales.
\end{problem}
\begin{problem}
    Sea $ABC$ un triangulo, y sean $D_1,D_2$ puntos en $BC$ tales que $AD_1,AD_2$ son isogonales respecto a $\angle BAC$. Demuestra que el circuncirculo de $ABC$ es tangente al de $AD_1D_2$.
\end{problem}

\begin{problem}
    Sea $ABC$ un triangulo, $P$ un punto en su interior, $Q$ su conjugado isogonal. Demuestra que $\angle BPC +\angle BQC =180+\angle A$.
\end{problem}
\subsection{Nivel II}
\begin{problem} %[Teorema de Steiner]
Sea $ABC$ un triangulo y sean $D$ y $E$ puntos en $BC$ tales que $AD$ y $AE$ son rectas isogonales. Entonces se cumple que: 
\[\frac{AB^2}{AC^2}=\frac{BD}{DC} \cdot \frac{BE}{EC}\]
\end{problem}

\begin{problem} %[Teorema de Pascal]
Sea $AEBDFC$ un hexágono cíclico, como se muestra. Supongamos que 
\[
P = \overline{AB} \cap \overline{DE}, \qquad
Q = \overline{CD} \cap \overline{FA}, \qquad
X = \overline{BC} \cap \overline{EF}.
\]
Demuestra que los puntos $P$, $X$ y $Q$ son colineales.
\end{problem}



\problemdiff{3}

\begin{problem}
     Sean $u$ y $v$ dos rectas que se intersectan en un punto $P$, y sea $T$ un punto del plano.  
 Sean $X$ y $Y$ los pies de altura del punto $T$ sobre las rectas $u$ y $v$, respectivamente.  
Demuestra que la recta $XY$ es perpendicular a la isogonal de la recta $PT$ con respecto a las rectas $u$ y $v$.

\end{problem}

\begin{problem}
    Sea $ABC$ un triangulo y $P$ un punto, sean $X,Y,Z$ las reflexiones de $P$ sobre $BC,CA,AB$. Sea $Q$ el circuncentro de $XYZ$. Demuestra que $P,Q$ son conjugados isogonales.
\end{problem}

\subsection{Nivel III}
\begin{problem}
Sean $P$ y $Q$ puntos conjugados isogonales en el triángulo $ABC$. 
Demuestra que los circuncentros de los triángulos $BPC$ y $BQC$ son inversos con 
respecto a la circunferencia circunscrita de $\triangle ABC$.
\end{problem}

\problemdiff{6}

\begin{problem} %[USAMO 2008]
     Sea $ABC$ un triángulo acutángulo y escaleno. Sean $M$, $N$ y $P$ los puntos medios de $\overline{BC}$, $\overline{CA}$ y $\overline{AB}$, respectivamente.  

Las mediatrices de $\overline{AB}$ y $\overline{AC}$ cortan al rayo $AM$ en los puntos $D$ y $E$, respectivamente. Sean $BD$ y $CE$ rectas que se intersectan en un punto $F$ dentro del triángulo $ABC$.  

Demuestra que los puntos $A$, $N$, $F$ y $P$ son concíclicos.
\end{problem}

\begin{problem} %[IMOSL 2012/G2]
    Sea $O$ el circuncentro y $H$ el ortocentro de un triángulo acutángulo $ABC$.  Demuestra que existen puntos $D$, $E$ y $F$ en los lados $BC$, $CA$ y $AB$, respectivamente, tales que
\[
OD + DH = OE + EH = OF + FH,
\]
y las rectas $AD$, $BE$ y $CF$ son concurrentes.
\end{problem}

\problemdiff{9}
\begin{problem} %[IMO SL 2000/G3]
Sea $O$ el circuncentro y $H$ el ortocentro de un triángulo acutángulo $ABC$. Demuestra que existen puntos $D$, $E$ y $F$ en los lados $BC$, $CA$ y $AB$, respectivamente, tales que
\[
OD + DH = OE + EH = OF + FH,
\]
y las rectas $AD$, $BE$ y $CF$ son concurrentes.
\end{problem}


\end{document}